\clearpage\thispagestyle{empty}
\pdfbookmark[0]{Quellen zum Brief vom 23. August 1807}{bode18070823-sources}
\begin{center}\textbf{Quellen zum Brief vom 23. August 1807}\end{center}
\noindent Diese KI-Übertragung in heutiges Deutsch folgt Gauß’ Brief an Johann Elert Bode vom 23. August 1807 in \href{https://gdz.sub.uni-goettingen.de/id/PPN236007467}{Gauß’ \textit{Werken}, Band~VI} (Göttingen, 1874), S.~284--285. Im \textit{Astronomischen Jahrbuch für das Jahr 1810} (Berlin, 1807) beginnt der Brief auf \href{https://www.digitale-sammlungen.de/en/view/bsb10538339?page=222}{S.~212}. Auf \href{https://www.digitale-sammlungen.de/en/view/bsb10538339?page=223}{S.~213} folgen die dritten Elemente. Die Ephemeride endet auf \href{https://www.digitale-sammlungen.de/en/view/bsb10538339?page=224}{S.~214}. Die Werkausgabe beschreibt die Ephemeride in eckigen Klammern. Die Digitalisate des Jahrbuchs stammen von der Bayerischen Staatsbibliothek.
\par\medskip
Gauß’ Abkürzung M.\,C. steht für die \textit{Monatliche Correspondenz}. Seine zweiten Elemente und seine Auskunft über ihre Übereinstimmung mit den Beobachtungen stehen auf \href{https://zs.thulb.uni-jena.de/rsc/viewer/jportal_derivate_00237836/Monatliche_Correspondenz_130168688_15_1807_0605.tif}{S.~598} in Franz Xaver von Zachs Bericht über Vesta vom Juni 1807, Beitrag XLIII, Band~XV, \href{https://zs.thulb.uni-jena.de/rsc/viewer/jportal_derivate_00237836/Monatliche_Correspondenz_130168688_15_1807_0597.tif}{S.~590}--600. Die Zeitschriftendigitalisate stammen von der Thüringer Universitäts- und Landesbibliothek Jena. Das editierbare Altarzeichen im Brief folgt seiner Kontur in der Werkausgabe. Übertragung und Satz wurden mit KI erstellt.
\clearpage
